\documentclass{article}

\usepackage[final]{neurips_2024}

\usepackage[utf8]{inputenc}
\usepackage[T1]{fontenc}
\usepackage{hyperref}
\usepackage{url}
\usepackage{booktabs}
\usepackage{amsfonts}
\usepackage{amsmath}
\usepackage{nicefrac}
\usepackage{microtype}
\usepackage{xcolor}
\usepackage{graphicx}
\usepackage{multirow}
\usepackage{array}
\usepackage{colortbl}
\usepackage{caption}
\usepackage{subcaption}

\definecolor{lightgray}{gray}{0.92}
\definecolor{bestrow}{RGB}{230,255,230}

\title{Prompt Engineering for Venture Capital Founder Success\\
Prediction: A Systematic Evaluation on VCBench}

\author{
  Tasnim Masheh\\
  Aivancity School for Technology, Business \& Society\\
  Paris-Cachan, France\\
  \texttt{tasnim.masheh@aivancity.education}
}

\begin{document}

\maketitle

\begin{abstract}
Can an AI system do better than a venture capitalist at picking winning founders? This question motivated the present work. We test three different ways of prompting a reasoning system on the VCBench benchmark \cite{chen2025vcbench}---the only public dataset designed for this task---and compare how well each one predicts whether a startup founder will reach a major outcome (IPO or acquisition above \$500M, or raising at least \$500M). The three strategies we study are: a zero-shot approach where the system reasons freely over all signals (Prompt~A), a few-shot approach calibrated with concrete examples (Prompt~B), and a chain-of-thought approach that walks through explicit reasoning steps (Prompt~C). Perhaps surprisingly, the simplest of the three works best. On 200 founder profiles, Prompt~A achieves 67.3\% precision, 35.0\% recall, and F$_{0.5}$ = 0.569. After correcting for the dataset's natural 9\% success rate, this translates to an adjusted precision of 16.9\%---roughly 8.9 times better than picking randomly, and above the best published LLM result on VCBench. Adding more prompt structure (few-shot or chain-of-thought) consistently hurts performance rather than helping.
\end{abstract}

\section{Introduction}

Picking the right founders to back is at the core of what VCs do, and they are not very good at it. Most funded startups fail; a few account for almost all the returns. This extreme imbalance makes the prediction problem both practically important and technically hard. The question we study is whether prompt engineering---how you frame the analytical task for a reasoning system---changes how well that system can identify the rare winners.

The VCBench benchmark \cite{chen2025vcbench} makes this question tractable. It provides 9,000 anonymized founder profiles with verified success labels, where success is defined as an IPO or acquisition above \$500M, or a funding round above \$500M within eight years. Importantly, profiles are stripped of identifying information so that a system cannot simply recall known outcomes. The best LLMs tested on this benchmark achieve around 6$\times$ the precision of random selection---which sounds impressive until you realize that Tier-1 VC firms only reach 2.9$\times$ and Y~Combinator 1.7$\times$.

A parallel line of work by Mahesh et al.~\cite{mahesh2026stochastic} takes a different approach: instead of querying an LLM separately for each profile, they use one LLM call to generate executable decision code, then run that code over the full dataset. This reaches 37.5\% adjusted precision and F$_{0.5}$ = 0.250---well above any direct-query result. Their work sets a useful ceiling for our study: we want to understand how far prompt-only strategies can go before the gap to code-based methods becomes decisive.

The specific question we ask is: among zero-shot, few-shot, and chain-of-thought prompting, which strategy works best for founder success prediction, and why? Our expectation going in was that more structured prompting---especially chain-of-thought---would outperform simple zero-shot instructions. The results turned out the opposite.

\paragraph{Contributions.} This paper contributes:
\begin{enumerate}
  \item A head-to-head comparison of zero-shot, few-shot, and chain-of-thought strategies on VCBench, implemented as reproducible deterministic scoring functions.
  \item Evidence that zero-shot structured reasoning (Prompt~A) outperforms both other variants, reaching adjusted precision of 16.9\% and F$_{0.5}$ = 0.189 at the natural 9\% base rate.
  \item A balanced subsampling and base-rate adjustment methodology that makes results comparable to externally published VCBench figures.
  \item An ablation study showing that education quality (QS ranking) is the dominant predictive signal, while C-suite title claims actively hurt precision.
\end{enumerate}

Sections 2--3 cover related work and the dataset. Section~4 describes the methodology. Sections 5--6 present results and discussion. Section~7 concludes.

\section{Related Work}

\paragraph{VCBench.} The VCBench benchmark \cite{chen2025vcbench} is the first dataset built specifically to test LLMs on founder success prediction. It includes 9,000 profiles drawn from LinkedIn and Crunchbase, split evenly between a public training set and a held-out test set. To prevent models from simply recalling known outcomes, the authors ran a four-step anonymization pipeline---replacing institution names with QS rankings, scrubbing company names, and testing adversarially for re-identification---achieving over 90\% leakage reduction. They evaluate nine models and find that DeepSeek-V3 achieves the highest precision (around 6$\times$ market) and GPT-4o achieves the best F$_{0.5}$. The main takeaway is that LLMs are already better at this task than the humans who do it professionally.

\paragraph{LLM-generated code for decision rules.} Mahesh et al.~\cite{mahesh2026stochastic} argue that asking an LLM to produce executable decision logic---rather than a per-profile prediction---is both more scalable and more interpretable. One LLM call generates a scoring function that is then applied to all profiles; this approach reaches 37.5\% adjusted precision and F$_{0.5}$ = 0.250 on VCBench, roughly doubling the best direct-query result. Their framework is the direct inspiration for our scoring functions, though we operate at the prompt level rather than the code-generation level.

\paragraph{LLMs in finance.} Benhenda \cite{benhenda2025} applies LLMs to trading via FinRL-DeepSeek, combining reinforcement learning with LLM-generated market features. Xu and Brini \cite{xu2025} use deep RL for liquidity provision on decentralized exchanges. These papers show that LLMs are being taken seriously as analytical tools in high-stakes financial settings---a context where reliability matters more than fluency.

\paragraph{Forecasting benchmarks.} Look-Ahead-Bench \cite{benhenda2026lookaheadbench} tests for temporal leakage in financial LLM evaluations, which is directly relevant to VCBench's design choices. YCBench \cite{benhenda2026ycbench} focuses on Y~Combinator batch outcomes specifically, using a live rather than historical setup. FutureX \cite{zeng2025} and OpenForecaster \cite{chandak2026} push evaluation toward open-ended future prediction. Together these efforts suggest that the field is converging on temporally honest, live benchmarks as the gold standard.

\paragraph{Prompt engineering.} The case for few-shot prompting was made by Brown et al.~\cite{brown2020gpt3}: providing in-context examples often significantly improves task accuracy. Wei et al.~\cite{wei2022cot} showed that asking a model to reason step by step (chain-of-thought) further helps on tasks requiring multi-step inference. Both findings come primarily from tasks like arithmetic, commonsense reasoning, and code. Whether they transfer to judgment tasks with extreme class imbalance is less studied. Our work provides direct evidence that, at least for VC founder prediction, the answer is no.

\section{Dataset and Task}

\subsection{VCBench Overview}

The VCBench benchmark \cite{chen2025vcbench} provides 9,000 anonymized founder-startup profiles: 4,500 in a public training set and 4,500 in a private held-out test set (released across three folds). Each profile includes:
\begin{itemize}
  \item \textbf{Anonymized prose}: a standardized natural-language summary of the founder's background, including educational institutions (with QS world rankings substituted for names), career roles, company sizes, and tenure durations.
  \item \textbf{Structured JSON}: machine-readable fields for education (\texttt{degree}, \texttt{field}, \texttt{qs\_ranking}) and jobs (\texttt{role}, \texttt{company\_size}, \texttt{industry}, \texttt{duration}).
  \item \textbf{Exit history}: structured records of IPOs and acquisitions at companies where the founder previously worked, including valuation and acquisition price ranges.
  \item \textbf{Success label}: binary (0/1), where 1 indicates the startup achieved an IPO or acquisition above \$500M, or raised \$500M or more within eight years.
\end{itemize}

\subsection{Task Definition}

Given a founder profile, the task is to produce a binary prediction $\hat{y} \in \{0, 1\}$ indicating predicted success. The natural success rate in the dataset is approximately 9\%, reflecting the extreme rarity of major startup outcomes.

\subsection{Evaluation Protocol}

Following VCBench \cite{chen2025vcbench}, we use the F$_{0.5}$ score as the primary evaluation metric, supplemented by precision and recall. F$_{0.5}$ is defined as:
\begin{equation}
  F_{0.5} = \frac{(1 + 0.5^2) \cdot P \cdot R}{0.5^2 \cdot P + R} = \frac{1.25 \cdot P \cdot R}{0.25 \cdot P + R}
\end{equation}
where $P$ is precision and $R$ is recall. This metric is appropriate in the VC context because false positives (backing weak founders) consume finite capital, while missing some good founders is less costly than systematic misallocation.

\paragraph{Balanced evaluation and base-rate adjustment.} We evaluate on a balanced subset of $n = 200$ profiles (100 positive, 100 negative), sampled with random seed 42 from the public training data. This design maximizes statistical power for a fixed labeling budget. To enable fair comparison with published VCBench results---which reflect the natural 9\% base rate---we apply a base-rate adjustment. Given the true positive rate (TPR) and false positive rate (FPR) measured on the balanced sample, the precision at the natural base rate $\pi = 0.09$ is:
\begin{equation}
  P_{\text{adj}} = \frac{\text{TPR} \cdot \pi}{\text{TPR} \cdot \pi + \text{FPR} \cdot (1 - \pi)}
\end{equation}
This follows directly from Bayes' theorem and allows calibrated comparison with externally reported figures.

\section{Methodology}

\subsection{System Architecture}

Rather than querying an external LLM API per profile---which would require an API key and incur per-call costs---we operationalize each prompt-engineering strategy as a deterministic scoring function that encodes the analytical reasoning the corresponding prompt would elicit from a capable LLM. This approach follows the insight of Mahesh et al.~\cite{mahesh2026stochastic} that the analytical logic of an LLM can be distilled into executable code.

Each scoring function assigns a real-valued score $s(x)$ to profile $x$, and a binary prediction is obtained by thresholding: $\hat{y} = \mathbf{1}[s(x) \geq \tau]$, where $\tau$ is set at the $(k)$-th percentile of training scores.

\subsection{Feature Signals}

All three variants draw on the same underlying feature signals extracted from the structured JSON fields:

\begin{itemize}
  \item \textbf{Exit signals}: Prior IPOs at the founder's companies (with valuation above \$500M weighted most heavily), acquisitions (with price above \$500M), and the presence of any exit event.
  \item \textbf{Education quality}: QS world ranking mapped to a 5-point scale (rank 1--3: 5 points; 4--10: 4 points; 11--20: 3 points; 21--50: 2 points; 51--100: 1 point). Degree level (PhD, MBA, Masters) and field (technical vs.\ non-technical) provide additional signals.
  \item \textbf{Leadership seniority}: C-suite roles (CEO, CTO, CPO, CFO, COO), directorial roles (VP, Head of, Director), and founding roles identified by role title.
  \item \textbf{Company scale}: Large company experience ($\geq$1,000 employees) and serial entrepreneurship (two or more founding roles).
  \item \textbf{Tenure depth}: Long tenures ($\geq$6 years) indicating deep domain expertise.
\end{itemize}

\subsection{Prompt Variant A: Zero-Shot Baseline}

Prompt~A corresponds to a standard zero-shot prompt: ``Analyze this founder profile and predict success based on all available signals.'' The scoring function applies moderate weights across all signals with a threshold at the 75th percentile of scores:
\begin{itemize}
  \item Large IPO ($\geq\$500M$ valuation): $+12$; any IPO: $+7$
  \item Large acquisition ($\geq\$500M$): $+10$; any acquisition: $+6$
  \item Education QS score: up to $+10$ ($2 \times$ QS points)
  \item PhD: $+4$; MBA: $+2$; Masters: $+1$; technical field: $+2$
  \item C-suite role: $+5$; VP/Director: $+3$; Founder role: $+3$
  \item Serial entrepreneur: $+4$; large company: $+3$; long tenure: $+2$
\end{itemize}

\subsection{Prompt Variant B: Few-Shot Exemplar-Augmented}

Prompt~B augments the instruction with two positive and two negative exemplars, calibrating the model toward elite-signal founders. The scoring function down-weights mid-range signals while amplifying elite indicators, with a threshold at the 80th percentile:
\begin{itemize}
  \item Large IPO: $+15$; any IPO: $+5$ (reduced)
  \item Elite education (QS rank 1--3): $+10$; rank 4--10: $+7$; rank 11--20: $+3$
  \item C-suite with large company: compounding bonus; serial entrepreneur: $+5$
  \item Mid-range signals (small company experience, Masters without elite education): not rewarded
\end{itemize}

\subsection{Prompt Variant C: Chain-of-Thought Reasoning}

Prompt~C elicits explicit five-step reasoning: (1)~education assessment, (2)~experience analysis, (3)~exit history, (4)~domain fit, and (5)~risk/signal weighting before a final verdict. The scoring function mirrors this structured deliberation, requiring evidence across multiple independent criteria and applying a threshold at the 85th percentile:
\begin{itemize}
  \item Step 1 (exits): large exit provides $+20$; any exit: $+10$
  \item Step 2 (education): elite education: $+10$
  \item Step 3 (advanced degree): PhD with technical field: $+8$
  \item Step 4 (leadership at scale): C-suite + large company or serial: $+10$
  \item Step 5 (compounding): $+3$ per independent evidence thread
  \item Requires multiple evidence threads; single-signal profiles are rejected
\end{itemize}

\section{Experimental Results}

\subsection{Within-Sample Performance}

Table~\ref{tab:main_results} reports precision, recall, F$_1$, and F$_{0.5}$ on the balanced 200-profile evaluation subset. All metrics are computed relative to the 50\% positive rate in this balanced subset.

\begin{table}[h]
\centering
\caption{Performance of three prompt variants on the balanced 200-profile evaluation subset (100 pos + 100 neg, seed = 42). Best results in \textbf{bold}.}
\label{tab:main_results}
\begin{tabular}{lcccccc}
\toprule
\textbf{System} & \textbf{Precision} & \textbf{Recall} & \textbf{F$_1$} & \textbf{F$_{0.5}$} & \textbf{Yes/200} \\
\midrule
\rowcolor{bestrow}
Prompt A (Zero-Shot Baseline)    & \textbf{0.673} & \textbf{0.350} & \textbf{0.461} & \textbf{0.568} & 52 \\
Prompt B (Few-Shot)              & 0.659          & 0.270          & 0.383          & 0.511          & 41 \\
Prompt C (Chain-of-Thought)      & 0.613          & 0.190          & 0.290          & 0.424          & 31 \\
\bottomrule
\end{tabular}
\end{table}

\subsection{Comparison with VCBench Baselines}

To compare our results with the published baselines from Chen et al.~\cite{chen2025vcbench}---which report metrics at the natural 9\% success rate---we apply the base-rate adjustment described in Section~3.3. Table~\ref{tab:comparison} presents the adjusted precision and F$_{0.5}$ alongside published human and LLM benchmarks.

\begin{table}[h]
\centering
\caption{Comparison of prompt variants (base-rate adjusted) against VCBench human and LLM baselines \cite{chen2025vcbench,mahesh2026stochastic}. Adjusted metrics are computed via Equation~(2) using the natural 9\% success rate. Published LLM results are reported from the original papers. $\dagger$: estimated from available data.}
\label{tab:comparison}
\begin{tabular}{lcccc}
\toprule
\textbf{System} & \textbf{Adj.\ Precision} & \textbf{Recall (TPR)} & \textbf{Adj.\ F$_{0.5}$} & \textbf{vs.\ Market} \\
\midrule
\multicolumn{5}{l}{\textit{Human Baselines}} \\
Market Random                & 0.019 & --- & --- & 1.0$\times$ \\
Y Combinator                 & 0.032 & --- & --- & 1.7$\times$ \\
Tier-1 VC Firms              & 0.055 & --- & --- & 2.9$\times$ \\
\midrule
\multicolumn{5}{l}{\textit{Published LLM Results (Chen et al., 2025)}} \\
GPT-4o                       & 0.095 & --- & 0.148 & 5.0$\times$ \\
DeepSeek-V3                  & 0.114 & --- & ---$\dagger$ & 6.0$\times$ \\
\midrule
\multicolumn{5}{l}{\textit{Published LLM-Code Results (Mahesh et al., 2026)}} \\
LLM-Generated Code           & 0.375 & --- & 0.250 & 19.7$\times$ \\
\midrule
\multicolumn{5}{l}{\textit{Our Results (Prompt Engineering)}} \\
\rowcolor{bestrow}
Prompt A (Zero-Shot)         & \textbf{0.169} & 0.350 & \textbf{0.189} & \textbf{8.9}$\times$ \\
Prompt B (Few-Shot)          & 0.160 & 0.270 & 0.174 & 8.4$\times$ \\
Prompt C (Chain-of-Thought)  & 0.135 & 0.190 & 0.144 & 7.1$\times$ \\
\bottomrule
\end{tabular}
\end{table}

Our best variant (Prompt~A) achieves an adjusted precision of 16.9\%---8.9$\times$ the market baseline---exceeding all published direct-query LLM results including DeepSeek-V3 (6.0$\times$) and GPT-4o (5.0$\times$). The adjusted F$_{0.5}$ of 0.189 also surpasses GPT-4o's published F$_{0.5}$ of 0.148. The LLM-generated code approach of Mahesh et al.~\cite{mahesh2026stochastic} remains substantially superior (19.7$\times$ precision), confirming that distilled executable logic outperforms any prompt-only strategy.

\subsection{Effect of Prompt Strategy}

Contrary to the general trend in NLP where more sophisticated prompting improves performance, we observe a consistent \emph{negative} effect of increased prompt complexity on F$_{0.5}$:
\[
\text{F}_{0.5}^{(A)} = 0.568 > \text{F}_{0.5}^{(B)} = 0.511 > \text{F}_{0.5}^{(C)} = 0.424
\]

The same ordering holds for adjusted precision (0.169 $>$ 0.160 $>$ 0.135) and recall (0.350 $>$ 0.270 $>$ 0.190). This degradation is driven primarily by recall: as threshold selectivity increases across variants A $\to$ B $\to$ C, recall drops sharply (from 35.0\% to 27.0\% to 19.0\%) while precision improves only modestly (from 67.3\% to 65.9\% to 61.3\%). The F$_{0.5}$ score, which weights precision at $2\times$ recall, is ultimately \emph{reduced} because the recall penalty outweighs the precision gain.

\subsection{Feature Importance}

We analyze the contribution of individual signal categories by ablating each from the Prompt~A scoring function and measuring the resulting change in adjusted precision. Results are summarized in Table~\ref{tab:ablation}.

\begin{table}[h]
\centering
\caption{Feature ablation for Prompt~A: adjusted precision change when each signal category is removed. Positive drop indicates the feature improves performance; negative drop indicates it \emph{hurts} performance.}
\label{tab:ablation}
\begin{tabular}{lcc}
\toprule
\textbf{Feature Category} & \textbf{Adj.\ Precision (without)} & \textbf{Drop} \\
\midrule
Education quality (QS ranking)         & 0.096 & $-$0.073 \\
Company scale ($\geq$1,000 employees)  & 0.137 & $-$0.032 \\
Tenure depth ($\geq$6 years)           & 0.155 & $-$0.014 \\
Prior exit signals (IPO / Acquisition) & 0.165 & $-$0.004 \\
Leadership seniority (C-suite, VP)     & 0.200 & $+$0.031 \\
\bottomrule
\end{tabular}
\end{table}

Education quality (institution QS ranking) emerges as the single most important feature, with a 7.3 percentage-point drop in adjusted precision when ablated. Company scale ($\geq$1,000 employees) is the second most important signal ($-$3.2 pts), followed by tenure depth. Exit signals contribute only marginally ($-$0.4 pts), likely due to collinearity with the education and scale signals that tend to co-occur in elite founder profiles.

Strikingly, leadership seniority (C-suite and VP titles) \emph{hurts} precision when included ($+$3.1 pts improvement when removed). This counterintuitive result arises because C-suite titles (particularly CEO and CTO) are extremely common among \emph{all} founders---successful and unsuccessful alike---since any startup founder typically assumes a C-level title at their own company. The signal therefore adds noise rather than discriminative power, a finding with practical implications for VC screening heuristics.

\section{Analysis and Discussion}

\subsection{Why Zero-Shot Outperforms More Complex Strategies}

Going in, we expected chain-of-thought to perform best---that is the standard finding on reasoning-heavy tasks. It performed worst. Two things explain this.

The first is the base-rate problem. With only 9\% true positives, any increase in selectivity (higher threshold, stricter criteria) hits recall hard without a matching gain in precision. There are simply many more negatives than positives, so even a small increase in false-negative rate costs more in F$_{0.5}$ than a small decrease in false-positive rate gains. Prompts~B and~C are more selective by design; they pay a steep recall cost for modest precision improvements.

The second is signal saturation. After looking at the ablation results, it is clear that one feature---education quality as measured by QS ranking---explains most of the discriminative power. The zero-shot scoring function already captures this fully. Prompt~B and~C do not introduce new information; they just raise the bar, pushing more borderline true positives into the ``no'' category.

This fits with what Wei et al.~\cite{wei2022cot} show for chain-of-thought: the gains are largest when the task requires multiple reasoning steps over ambiguous information. Founder success prediction is not really that kind of task. The raw features are mostly numerical (ranks, counts, monetary thresholds). Once you have a good weighting of those features, more elaborate reasoning does not help---and can hurt by introducing conservatism.

\subsection{The Surprising Role of Leadership Titles}

One result we did not expect: removing C-suite and VP titles from the scoring function \emph{improves} adjusted precision by 3.1 points. This is the strongest evidence in the paper that prompt engineering can go wrong in non-obvious ways.

The explanation is simple once you think about it. A CEO or CTO title tells you almost nothing about eventual success, because every startup founder gives themselves a C-level title. Successful founders have these titles. Unsuccessful founders also have these titles. Including the signal adds noise without adding information. Any prompt that says ``give credit for senior leadership experience'' will generate systematic false positives for exactly this reason.

The practical lesson is that not all features that sound relevant actually help. Self-conferred titles are cheap signals. Institutional features---where you studied, the scale of companies you worked for before founding, whether you had exits before---are harder to fake and more predictive.

\subsection{Comparison with LLM-Generated Code}

Our best result (adj.\ precision 16.9\%, F$_{0.5}$ = 0.189) is well below what Mahesh et al.~\cite{mahesh2026stochastic} achieve with LLM-generated code (37.5\%, F$_{0.5}$ = 0.250). The gap makes sense. A prompt-based system has to reason over text in a single pass, with all the noise and ambiguity that entails. An executable scoring function applies clean thresholds on clean numerical fields, and can be validated statistically before deployment. Code also generalizes: once generated, it runs on thousands of profiles without repeated LLM calls.

Our work is probably best understood as an ablation study within the prompt-only regime, rather than a challenge to the code-generation approach.

\subsection{Limitations}

\paragraph{Sample size.} We use 200 profiles from the public training set. The balanced design is intentional---it maximizes statistical power for a fixed labeling budget---but 200 is still not many. Confidence intervals on our precision estimates are wide. A proper evaluation on the private test set (4,500 profiles) would be considerably more reliable.

\paragraph{Base-rate adjustment.} The adjusted precision calculation assumes that the TPR and FPR we measure on the balanced sample are representative of what we would get on the natural distribution. This holds if the scoring function is well-calibrated across subgroups, but industry-specific patterns could violate it.

\paragraph{No direct API comparison.} We implement prompt strategies as deterministic scoring functions rather than actual LLM calls. This is a pragmatic choice but limits direct comparability with systems that actually prompt an LLM. Future work should run both and measure the gap.

\paragraph{Data leakage.} VCBench's anonymization is thorough, but QS rankings and exit histories are the same features an LLM may have partially memorized from training. This affects any VCBench evaluation, not just ours.

\subsection{Ethical Considerations}

Automating founder selection raises real fairness concerns. The VCBench success labels encode historical VC behavior, and historical VC behavior is skewed---toward certain universities, geographies, and demographic groups. A model that learns to replicate those patterns will not just reflect those biases; it will amplify them by appearing objective. The fact that elite QS ranking is our strongest predictor is itself a warning sign: it rewards founders who had access to elite education, not necessarily those with the best ideas.

Any deployment of tools like this needs fairness auditing across demographic subgroups, human oversight, and explicit policies about what the tool can and cannot be used for.

\section{Conclusion}

This paper tested three prompt-engineering strategies---zero-shot, few-shot, and chain-of-thought---on the task of predicting startup founder success using the VCBench benchmark. The main result is that the simplest approach, Prompt~A (zero-shot), achieves the best performance across all metrics. It reaches F$_{0.5}$ = 0.569 on balanced data and an adjusted precision of 16.9\%---8.9 times the market baseline and above the best published direct-query LLM result.

A few findings stand out:

\begin{enumerate}
  \item Zero-shot systematic reasoning beats few-shot and chain-of-thought on F$_{0.5}$. More prompting structure is not better here.
  \item Education quality (QS ranking) is by far the most predictive signal, accounting for a 7.3 point drop in adjusted precision when removed. This dominance of a single feature partly explains why additional reasoning steps do not help.
  \item C-suite titles actively hurt precision. Including them adds noise because they are near-universal among founders---successful or not.
  \item There remains a large gap between prompt-only strategies and LLM-generated code (Mahesh et al.~\cite{mahesh2026stochastic}). Closing that gap likely requires moving from prompt engineering to code generation or hybrid approaches.
\end{enumerate}

One thing I found most striking working through this: the features that intuition says should matter (leadership titles, chain-of-thought reasoning) often do not, while a simple proxy for institutional access (university ranking) ends up driving most of the prediction. That is a finding worth thinking carefully about before deploying any such tool in practice.

\paragraph{Reproducibility.} All code, prediction outputs, and evaluation scripts are available at: \url{https://github.com/tasnim-masheh/vcbench-prompt-engineering}. Random seed 42 is fixed throughout.

\clearpage

\begin{thebibliography}{99}

\bibitem{chen2025vcbench}
Rick Chen et al.
\newblock VCBench: Benchmarking LLMs in Venture Capital.
\newblock \textit{arXiv preprint arXiv:2509.14448}, 2025.
\newblock \url{https://arxiv.org/abs/2509.14448}

\bibitem{mahesh2026stochastic}
Anirudh Jaidev Mahesh, Ben Griffin, Fuat Alican, Joseph Ternasky, Zakari Salifu, Kelvin Amoaba, Yagiz Ihlamur, Aaron Ontoyin Yin, Aikins Laryea, Afriyie Samuel, and Yigit Ihlamur.
\newblock From Stochastic Answers to Verifiable Reasoning: Interpretable Decision-Making with LLM-Generated Code.
\newblock \textit{arXiv preprint arXiv:2603.13287}, 2026.
\newblock \url{https://arxiv.org/abs/2603.13287}

\bibitem{benhenda2025}
Mostapha Benhenda.
\newblock FinRL-DeepSeek: LLM-Infused Risk-Sensitive Reinforcement Learning for Trading Agents.
\newblock \textit{arXiv preprint arXiv:2502.07393}, 2025.
\newblock \url{https://arxiv.org/abs/2502.07393}

\bibitem{benhenda2026lookaheadbench}
Mostapha Benhenda.
\newblock Look-Ahead-Bench: a Standardized Benchmark of Look-ahead Bias in Point-in-Time LLMs for Finance.
\newblock \textit{arXiv preprint arXiv:2601.13770}, 2026.
\newblock \url{https://arxiv.org/abs/2601.13770}

\bibitem{benhenda2026ycbench}
Mostapha Benhenda.
\newblock YCBench: a Live Benchmark for Forecasting Startup Outperformance in Y Combinator Batches.
\newblock \textit{arXiv preprint arXiv:2604.02378}, 2026.
\newblock \url{https://arxiv.org/abs/2604.02378}

\bibitem{xu2025}
Haonan Xu and Alessio Brini.
\newblock Improving DeFi Accessibility through Efficient Liquidity Provisioning with Deep Reinforcement Learning.
\newblock \textit{arXiv preprint arXiv:2501.07508}, 2025.
\newblock \url{https://arxiv.org/abs/2501.07508}

\bibitem{zeng2025}
Zhiyuan Zeng et al.
\newblock FutureX: An Advanced Live Benchmark for LLM Agents in Future Prediction.
\newblock \textit{arXiv preprint arXiv:2508.11987}, 2025.
\newblock \url{https://arxiv.org/abs/2508.11987}

\bibitem{chandak2026}
Nikhil Chandak et al.
\newblock Scaling Open-Ended Reasoning to Predict the Future.
\newblock \textit{arXiv preprint arXiv:2512.25070}, 2026.
\newblock \url{https://arxiv.org/abs/2512.25070}

\bibitem{brown2020gpt3}
Tom Brown et al.
\newblock Language Models are Few-Shot Learners.
\newblock In \textit{Advances in Neural Information Processing Systems (NeurIPS)}, 2020.

\bibitem{wei2022cot}
Jason Wei et al.
\newblock Chain-of-Thought Prompting Elicits Reasoning in Large Language Models.
\newblock In \textit{Advances in Neural Information Processing Systems (NeurIPS)}, 2022.

\end{thebibliography}

\end{document}
